\section{Metric geometry and the triangle portfolio}
\label{sec:triangles}

This chapter gives finite algebraic proofs of the familiar metric and triangle
theorems.  Coordinates are the certified exact scalar codes of the foundational
chapter.  Ordered-field notation is an abbreviation for computations on those
codes, not an additional higher-order axiom.  In particular, a square root below
means its nonnegative, certified algebraic root.  Division is used only after a
nonzero denominator has been exhibited.  The coordinate calculations give
\ipra{} proof schemata by the certificate interpretation established there.
They do not assert a decision procedure for equality of arbitrary Cauchy reals.

The classical statements covered here include the triangle propositions in
Euclid's Books I, II and VI; our proofs use coordinates and explicitly account
for their constructive data.  For example, the original statements of
Pythagoras and the angle-bisector theorem are in I.47 and VI.3,
respectively.\footnote{Euclid, \emph{Elements}, I.47 and VI.3, in David Joyce's
online edition:
\url{https://aleph0.clarku.edu/~djoyce/java/elements/bookI/propI47.html} and
\url{https://aleph0.clarku.edu/~djoyce/java/elements/bookVI/propVI3.html}.
These references identify the classical statements; the arithmetic proofs
below are supplied in full.}

\subsection{Definitions and two-coordinate certificates}

For vectors $u=(u_1,u_2)$ and $v=(v_1,v_2)$ put
\[
 u\mathbin{\cdot}v=u_1v_1+u_2v_2,\qquad
 [u,v]=u_1v_2-u_2v_1,\qquad
 |u|^2=u\mathbin{\cdot}u,\qquad
 Ju=(-u_2,u_1).
\]
The notation $|u|$ denotes the nonnegative square root of $|u|^2$.
Thus $|P-Q|$ is a distance.  A noncollinear triangle $(A,B,C)$ has
$[B-A,C-A]\ne0$.  Unless orientation matters, write
\[
 a=|B-C|,\quad b=|C-A|,\quad c=|A-B|,\quad
 D=|[B-A,C-A]|,\quad \Delta=D/2,\quad s=(a+b+c)/2.
\]
Its side lengths and $D$ are strictly positive.  Reflection, or an interchange
of $B,C$, permits us to calculate with $A=0$, $B=u$, $C=v$ and $[u,v]=D>0$.
This is an explicit finite coordinate change.  It is not an assumption that
an arbitrary real sign can be decided.

\begin{lemma}[Gram identity and quarter-turn identities]
\label{tri:gram}
For all $u,v$,
\[
 (u\mathbin{\cdot}v)^2+[u,v]^2=|u|^2|v|^2,\qquad
 (Ju)\mathbin{\cdot}v=[u,v],\qquad
 u\mathbin{\cdot}(Jv)=-[u,v],\qquad |Ju|^2=|u|^2.
\]
If $u\ne0$, every $v$ has the unique decomposition
\[
 v=\frac{u\mathbin{\cdot}v}{|u|^2}u+
   \frac{[u,v]}{|u|^2}Ju.
\]
\end{lemma}
\begin{proof}
The first left-hand side expands to
$u_1^2v_1^2+2u_1u_2v_1v_2+u_2^2v_2^2+
u_1^2v_2^2-2u_1u_2v_1v_2+u_2^2v_1^2$.
Cancellation leaves $(u_1^2+u_2^2)(v_1^2+v_2^2)$.
The three quarter-turn equalities follow by substituting $Ju=(-u_2,u_1)$.
For the decomposition, multiply through by $|u|^2$ and expand its two
coordinates.  They become $|u|^2v_1$ and $|u|^2v_2$.
Taking the dot product with $u$ and with $Ju$ recovers the two coefficients,
which proves uniqueness.  The denominator is positive because a sum of two
squares can vanish only if both coordinates vanish.
\end{proof}

\begin{theorem}[Cauchy--Schwarz, metric axioms, and strict triangle inequality]
\label{tri:metric}
For all vectors $u,v$,
\[
 -|u||v|\le u\mathbin{\cdot}v\le |u||v|,
 \qquad |u+v|\le |u|+|v|.
\]
Distance is symmetric, nonnegative, and zero precisely on equal points.
If $[u,v]\ne0$, then $|u+v|<|u|+|v|$.
Consequently a noncollinear triangle satisfies all three strict inequalities
$a<b+c$, $b<c+a$, and $c<a+b$.
\end{theorem}
\begin{proof}
The Gram identity gives $(u\mathbin{\cdot}v)^2\le |u|^2|v|^2$.
For a nonnegative $q$, the inequalities $p>q$ or $p<-q$ would imply
$p^2>q^2$, by factoring the difference of squares and multiplying positive
quantities.  Apply this with $p=u\mathbin{\cdot}v$ and $q=|u||v|$.
Now
\[
 |u+v|^2=|u|^2+2u\mathbin{\cdot}v+|v|^2
 \le (|u|+|v|)^2.
\]
Both sides being squares of nonnegative quantities, the same factorization
gives the claimed inequality for their square roots.  If $[u,v]\ne0$,
the Gram inequality is strict, so the last inequality is strict as well.
Symmetry and the zero criterion follow directly from the sum-of-squares
definition.  Apply the strict inequality to the two edge vectors along a
broken path between each pair of vertices.  Their determinants differ from
$[B-A,C-A]$ only by sign, and are therefore nonzero.
\end{proof}

\begin{proposition}[Explicit Euclidean motions and similarities]
\label{tri:motions}
Let $e,e'$ be unit vectors and let $\varepsilon^2=1$.  The map
\[
 Qx=(x\mathbin{\cdot}e)e'+
       \varepsilon(x\mathbin{\cdot}Je)Je'
\]
preserves dot products and distances.  It has determinant $\varepsilon$.
For $\lambda>0$, the map $x\mapsto P+\lambda Q(x-A)$ multiplies all
distances by $\lambda$, preserves all unoriented angles, and multiplies
unsigned triangle areas by $\lambda^2$.
\end{proposition}
\begin{proof}
Write $x=\alpha e+\beta Je$ and $y=\gamma e+\delta Je$ using
Lemma~\ref{tri:gram}.  Orthogonality and the unit equations give
$x\mathbin{\cdot}y=\alpha\gamma+\beta\delta$ and
$Qx\mathbin{\cdot}Qy=\alpha\gamma+\varepsilon^2\beta\delta$.
The two are equal.  The same computation with $x-y$ proves distance
preservation.  The determinant of the matrix with columns $e,Je$ is one,
as is that with columns $e',Je'$; the intervening diagonal matrix has
determinant $\varepsilon$.  Scaling introduces a factor $\lambda^2$
in dot products and determinants.  The ratios defining angles below
therefore remain unchanged, and their determinant absolute values scale
by $\lambda^2$ as asserted.
\end{proof}

\subsection{Angles without transcendental functions}

For nonzero vectors $u,v$, the unoriented angle between their rays is the pair
\[
 \operatorname{Ang}(u,v)=
 \left(\frac{u\mathbin{\cdot}v}{|u||v|},
       \frac{|[u,v]|}{|u||v|}\right).
\]
It belongs to the upper unit semicircle.  Its components will be called
$\cos\alpha$ and $\sin\alpha$; this is an algebraic definition, not an appeal
to already constructed analytic sine and cosine.  Equality of angles means
equality of these pairs.  A right angle is $(0,1)$; a straight angle is
$(-1,0)$.  An angle is acute if its first coordinate is positive and obtuse if
that coordinate is negative.  For the open upper semicircle, angular order
is circular order: $z<w$ means $[z,w]>0$.  This also agrees with decreasing
first coordinate.  Indeed, for $z=(c,t)$ and $w=(d,r)$ there,
\[
 (t+r)[z,w]=(c-d)(1+z\mathbin{\cdot}w),
\]
and both $t+r$ and $1+z\mathbin{\cdot}w$ are positive.  The latter is
$|z+w|^2/2$, and $z+w$ has positive second coordinate.

Addition of direction angles uses the pair product
\[
 (c,t)*(d,r)=(cd-tr,cr+td).
\]
If two open upper-semicircle angles have their product in the open upper
semicircle, this product is their sum with no passage through a straight
angle: $[z,z*w]=r>0$ and $[w,z*w]=t>0$.  More generally, half-plane signs
record the finite carry when a sum crosses a half-turn.  These observations
are a precise algebraic meaning for the angle sums used below.  Supplement
is $(c,t)\mapsto(-c,t)$.

\begin{theorem}[Pythagoras and its converse]
\label{tri:pythagoras}
For three distinct vertices, the angle at $A$ is right if and only if
$a^2=b^2+c^2$.
\end{theorem}
\begin{proof}
Writing $u=B-A$, $v=C-A$ gives
\[
 a^2=|v-u|^2=|u|^2+|v|^2-2u\mathbin{\cdot}v
             =c^2+b^2-2u\mathbin{\cdot}v.
\]
Thus the asserted equality is equivalent to $u\mathbin{\cdot}v=0$.
The Gram identity then gives $|[u,v]|=|u||v|$, so the angle pair is
$(0,1)$.  Conversely that pair has zero dot product.  All cancellation is
by the nonzero constant $2$ and the positive side lengths.
\end{proof}

\begin{theorem}[Law of cosines]
\label{tri:cosine}
For a noncollinear triangle with interior angle pairs $\alpha,\beta,\gamma$,
\[
 a^2=b^2+c^2-2bc\cos\alpha,
\quad b^2=c^2+a^2-2ca\cos\beta,
\quad c^2=a^2+b^2-2ab\cos\gamma.
\]
\end{theorem}
\begin{proof}
The expansion in the preceding proof and the definition
$bc\cos\alpha=(B-A)\mathbin{\cdot}(C-A)$ give the first identity.
Perform the same expansion with origins $B$ and $C$ for the other two.
\end{proof}

\begin{theorem}[Triangle angle sum and exterior-angle theorem]
\label{tri:angle-sum}
The three interior angles of a noncollinear triangle sum to a straight
angle.  Each exterior angle is the sum of the two remote interior angles
and is greater than either of them.
\end{theorem}
\begin{proof}
Use positive orientation $A=0$, $B=u$, $C=v$ and define unit vectors
\[
 e=u/c,\qquad f=v/b,\qquad g=(v-u)/a.
\]
Regard a pair $(x,y)$ as a formal complex number, with multiplication $*$
above and conjugation $(x,y)\mapsto(x,-y)$.  The three interior angle pairs
are exactly
\[
 z_A=\overline e f,\qquad z_B=\overline g(-e),\qquad
 z_C=\overline f g.
\]
Their imaginary parts are $D/(bc)$, $D/(ac)$ and $D/(ab)$, all positive.
Associativity, commutativity and $\overline e e=\overline f f=
\overline g g=1$ give
\[
 z_Az_B=-\overline{z_C}=(-\cos\gamma,\sin\gamma).
\]
The product is in the open upper semicircle.  The angle-addition rule just
defined therefore makes it the sum of $\alpha$ and $\beta$ without a
half-turn carry.  It is the supplement of $\gamma$, so adjoining
$\gamma$ gives a straight angle.  The supplement is the exterior angle at
$C$.  Its determinants with $z_A$ and $z_B$ are respectively
$\sin\beta>0$ and $\sin\alpha>0$, so it is greater than each.  Cycling the
vertices proves the other two exterior-angle statements.
\end{proof}

\begin{theorem}[At most one nonacute angle; at least two acute angles]
\label{tri:acute}
A noncollinear triangle cannot have two right or obtuse angles.
For certified exact coordinates, at least two of its three angles are acute.
\end{theorem}
\begin{proof}
Suppose the angle at $A$ is not acute.  With $u=B-A$, $v=C-A$, put
$p=u\mathbin{\cdot}v\le0$.  The dot products defining the angles at $B,C$
are $|u|^2-p>0$ and $|v|^2-p>0$.  They are acute.  The same argument applies
at any designated nonacute vertex.  This proves the first assertion
without any case choice.  For the second, the finite scalar comparison
procedure either certifies all three dot products positive or supplies one
that is nonpositive; in the latter case the calculation just made certifies
the other two positive.  This final decidable split is asserted only for the
exact coordinate representation, not for arbitrary Cauchy names.
\end{proof}

\begin{theorem}[Area--sine identity and sine law]
\label{tri:sine}
For every noncollinear triangle,
\[
 2\Delta=bc\sin\alpha=ca\sin\beta=ab\sin\gamma,
 \qquad
 \frac a{\sin\alpha}=\frac b{\sin\beta}
                      =\frac c{\sin\gamma}=\frac{abc}{2\Delta}.
\]
\end{theorem}
\begin{proof}
The first three equalities are precisely the definitions of the three
positive sine components: the determinants at all vertices have absolute
value $D=2\Delta$.  Divide by $abc>0$ and then by each positive sine to
obtain the displayed ratios.  No inverse trigonometric function or limit is
involved.
\end{proof}

\begin{theorem}[Side--angle ordering]
\label{tri:side-angle}
In a noncollinear triangle, $a>b$ if and only if $\alpha>\beta$;
$a=b$ if and only if $\alpha=\beta$.  The analogous statements hold for
the other pairs of sides and opposite angles.
\end{theorem}
\begin{proof}
Subtract the cosine-law expressions to get the exact factorization
\[
 \cos\alpha-\cos\beta
 =\frac{(a-b)\bigl(c^2-(a+b)^2\bigr)}{2abc}.
\]
By strict triangle inequality, $c<a+b$, so the second numerator factor is
negative.  The denominator is positive.  Thus $a>b$ is equivalent to
$\cos\alpha<\cos\beta$, and the upper-semicircle order identity above
identifies this with $\alpha>\beta$.  If $a=b$, the cosine components are
equal, and the sine law makes the positive sine components equal too.
Conversely equal angle pairs have equal sines, so the sine law gives $a=b$.
\end{proof}

\subsection{Congruence, similarity, and constructions}

Triangles are \emph{congruent} here when one explicit Euclidean motion of
Proposition~\ref{tri:motions}, including a possible reflection, sends their
ordered vertices to one another.  They are \emph{similar} if a positive
scaling followed by such a motion does so.  These are witness-bearing
definitions, stronger than a mere list of invariant equalities.

\begin{theorem}[SSS congruence, with an explicit motion]
\label{tri:sss}
Two noncollinear triangles with the same three corresponding side lengths
are congruent.
\end{theorem}
\begin{proof}
Translate their first vertices to zero, and call the remaining vectors
$u,v$ and $u',v'$.  Set $c=|u|=|u'|>0$, $e=u/c$, $e'=u'/c$.
The cosine expansion gives
\[
 p=u\mathbin{\cdot}v=u'\mathbin{\cdot}v'
   =\frac{|u|^2+|v|^2-|v-u|^2}{2}.
\]
Let $d=[u,v]\ne0$ and $d'=[u',v']\ne0$.  The Gram identity and equality
of corresponding lengths give $d'^2=d^2$.  Thus
$\varepsilon=d'/d$ satisfies $\varepsilon^2=1$; no sign choice is needed.
The frame decompositions are
\[
 u=ce,\quad v=(p/c)e+(d/c)Je,\qquad
 u'=ce',\quad v'=(p/c)e'+(d'/c)Je'.
\]
The motion $Q$ in Proposition~\ref{tri:motions}, with this $\varepsilon$,
sends $u$ to $u'$ and $v$ to $v'$.  Restoring the two translations gives
the required witness.
\end{proof}

\begin{theorem}[SAS congruence]
\label{tri:sas}
Equality of two corresponding side lengths and of their included angle
implies congruence of noncollinear triangles.
\end{theorem}
\begin{proof}
If $b,c$ and $\alpha$ agree, the cosine law gives equality of $a^2$.
The positive square roots $a$ are therefore equal.  Apply the explicit SSS
construction.  This proves the assertion for any included vertex after
renaming the vertices.
\end{proof}

\begin{theorem}[ASA and AAS congruence]
\label{tri:asa}
Equality of two corresponding angles and of one corresponding side implies
congruence of noncollinear triangles.  In particular both ASA and AAS are
valid.
\end{theorem}
\begin{proof}
The angle-sum identity determines the third angle from the other two, so
all three pairs agree.  Suppose the common side is $c$.  The positive
$\sin\gamma$ is common, and the sine law gives the explicit formulas
\[
 a=c\frac{\sin\alpha}{\sin\gamma},\qquad
 b=c\frac{\sin\beta}{\sin\gamma}.
\]
Thus all side lengths agree, and SSS supplies a motion.  If another side
was specified, use the corresponding cyclic formulas.  In each case the
only new denominator is the sine of a nondegenerate interior angle.
\end{proof}

\begin{theorem}[AAA, SSS and SAS similarity]
\label{tri:similarity}
Two noncollinear triangles are similar if their corresponding angles agree;
if their corresponding sides are in one positive ratio; or if two
corresponding sides are in one positive ratio and their included angles
agree.  Similarity preserves all angle pairs and makes all corresponding
side ratios equal.
\end{theorem}
\begin{proof}
For AAA, let $\lambda=c'/c>0$ and scale the first triangle by $\lambda$.
Its side $c$ then agrees with $c'$ and its angles are unchanged, so ASA
gives a congruence to the second triangle.  For proportional SSS, the
same scaling makes all three sides agree, so use SSS.  For proportional
SAS, scaling makes the designated two sides agree, so use SAS.
Conversely Proposition~\ref{tri:motions} computes exactly how a similarity
changes lengths and angle pairs.  These constructions exhibit the scale
and the motion rather than merely asserting their existence.
\end{proof}

\begin{theorem}[Pons asinorum and the equilateral angle theorem]
\label{tri:isosceles}
In a noncollinear triangle, $AB=AC$ if and only if the base angles at
$B,C$ are equal.  Their exterior supplements are equal as well.
An equilateral triangle has all three angle pairs
$(1/2,\sqrt3/2)$, the algebraic representation of $60$ degrees.
\end{theorem}
\begin{proof}
The first equivalence is the equality case of
Theorem~\ref{tri:side-angle}; applying the supplement operation gives the
exterior statement.  If $a=b=c=\ell>0$, the cosine law gives
$\cos\alpha=(\ell^2+\ell^2-\ell^2)/(2\ell^2)=1/2$ at every vertex.
The positive sine has square $1-1/4=3/4$, hence is $\sqrt3/2$.
This pair has cube $(-1,0)$, and three successive positive equal turns
fill one straight angle by Theorem~\ref{tri:angle-sum}; this is exactly
one third of $180$ degrees, independently of a radian construction.
\end{proof}

\begin{proposition}[A rational SSA counterexample]
\label{tri:ssa}
Two sides and a nonincluded angle do not determine a triangle up to
congruence, even when all vertex coordinates are rational.
\end{proposition}
\begin{proof}
Take $A=(0,0)$, $B=(5,0)$ and
\[
 C_1=(6/5,8/5),\qquad C_2=(12/5,16/5).
\]
In both triangles $AB=5$, $BC_i^2=17$, and the angle at $A$ is the pair
$(3/5,4/5)$.  Direct sums of squares give
$AC_1=2$ and $AC_2=4$.  A congruence would preserve that distance, so none
exists.  Both determinants $[B,C_i]$ are positive.  Thus the example does
not exploit a degenerate triangle or an undecidable branch.
\end{proof}

\begin{theorem}[Equilateral triangles on a given segment]
\label{tri:equilateral-construction}
If $A\ne B$ and $d=B-A$, precisely the two points
\[
 C_\pm=\frac{A+B}{2}\ \pm\ \frac{\sqrt3}{2}Jd
\]
form equilateral triangles with side $AB$.  Their determinant signs are
opposite, so one lies on each side of the supporting line.
\end{theorem}
\begin{proof}
The vectors from $A,B$ to $C_\pm$ are
$d/2\pm(\sqrt3/2)Jd$ and $-d/2\pm(\sqrt3/2)Jd$.
Orthogonality gives their squared lengths as
$(1/4+3/4)|d|^2=|d|^2$.
Also $[d,C_\pm-A]=\pm(\sqrt3/2)|d|^2\ne0$.
Conversely let $C$ have the required distances and put
$w=C-(A+B)/2$.  Subtracting the two distance equations gives
$w\mathbin{\cdot}d=0$.  Lemma~\ref{tri:gram} therefore gives $w=tJd$.
One distance equation becomes $(1/4+t^2)|d|^2=|d|^2$, hence
$t^2=3/4$.  Factoring
$(t-\sqrt3/2)(t+\sqrt3/2)=0$ and using the exact scalar comparison
procedure gives exactly the two displayed possibilities.
\end{proof}

\subsection{Area, medians, and right-triangle projections}

\begin{theorem}[Heron's formula]
\label{tri:heron}
Every noncollinear triangle satisfies
\[
 16\Delta^2=2a^2b^2+2b^2c^2+2c^2a^2-a^4-b^4-c^4
           =(a+b+c)(-a+b+c)(a-b+c)(a+b-c),
\]
and therefore $\Delta^2=s(s-a)(s-b)(s-c)$.
\end{theorem}
\begin{proof}
With $p=u\mathbin{\cdot}v=(b^2+c^2-a^2)/2$, the Gram identity yields
\[
 16\Delta^2=4D^2=4b^2c^2-(b^2+c^2-a^2)^2.
\]
Expanding the square gives the first polynomial expression.  Factoring
as a difference of squares gives
\[
 \bigl(2bc-(b^2+c^2-a^2)\bigr)
 \bigl(2bc+(b^2+c^2-a^2)\bigr)
 =(a^2-(b-c)^2)((b+c)^2-a^2),
\]
which factors into the four displayed linear factors.  Replacing them
by $2s$, $2(s-a)$, $2(s-b)$ and $2(s-c)$ and cancelling $16$ proves the
last identity.  All four factors are positive by the strict triangle
inequalities, so this is also a certified positive square-root formula
for the area.
\end{proof}

\begin{theorem}[Constructing a triangle from its side lengths]
\label{tri:side-construction}
For positive certified side lengths $a,b,c$, a noncollinear triangle with
those lengths exists if and only if the three strict triangle inequalities
hold.  A witness is
\[
 A=(0,0),\qquad B=(c,0),\qquad C=(x,y),\qquad
 x=\frac{b^2+c^2-a^2}{2c},\quad y=\sqrt{b^2-x^2}>0.
\]
\end{theorem}
\begin{proof}
Necessity was proved in Theorem~\ref{tri:metric}.  For sufficiency, direct
factorization gives
\[
 4c^2(b^2-x^2)
 =(a+b+c)(-a+b+c)(a-b+c)(a+b-c)>0.
\]
Thus the nonnegative root $y$ is strictly positive.  The definitions give
$AC^2=x^2+y^2=b^2$ and
$BC^2=(x-c)^2+y^2=b^2-2cx+c^2=a^2$.
The determinant is $cy>0$.  The three distances are the designated
positive roots.  These formulas use a single positive square-root
construction and provide all denominator and nondegeneracy witnesses.
\end{proof}

\begin{theorem}[Parallelogram law and Apollonius's median theorem]
\label{tri:apollonius}
For all vectors,
\[
 |u+v|^2+|u-v|^2=2|u|^2+2|v|^2.
\]
If $M=(B+C)/2$ is the midpoint of a triangle side, then
\[
 AM^2=\frac{2AB^2+2AC^2-BC^2}{4}.
\]
\end{theorem}
\begin{proof}
Expand the two squares in the first expression.  Their dot-product
terms are $2u\mathbin{\cdot}v$ and $-2u\mathbin{\cdot}v$, and cancel.
Apply the result to $u=B-A$, $v=C-A$ and use
$M-A=(u+v)/2$ and $B-C=u-v$.  Rearranging gives the second formula.
\end{proof}

\begin{theorem}[Stewart's theorem]
\label{tri:stewart}
Let $T$ lie on $BC$, let $m=BT$, $n=TC$, and suppose $m,n\ge0$ and
$a=m+n>0$.  With $d=AT$,
\[
 b^2m+c^2n=a(d^2+mn).
\]
\end{theorem}
\begin{proof}
The certified betweenness data give
$T=(nB+mC)/a$.  For $u=B-A$, $v=C-A$ this means
$a(T-A)=nu+mv$.  Consequently
\[
 a^2d^2=n^2c^2+m^2b^2+2mn(u\mathbin{\cdot}v).
\]
Substitute $2u\mathbin{\cdot}v=b^2+c^2-a^2$ and collect terms.  The
right side becomes
$na c^2+ma b^2-mna^2$ because $m+n=a$.
Cancel the positive $a$ and rearrange.  Endpoints $m=0$ or $n=0$ are
included, so this proof does not require an unmentioned interior-point
hypothesis.
\end{proof}

\begin{theorem}[Right-triangle altitude and projection theorems]
\label{tri:right-altitude}
Suppose $\angle A$ is right, and let $F$ be the perpendicular projection
of $A$ on $BC$.  Then $F$ is strictly between $B,C$ and
\[
 BF=\frac{c^2}{a},\qquad FC=\frac{b^2}{a},\qquad
 AF=\frac{bc}{a},\qquad AF^2=BF\,FC.
\]
In particular $AB^2=BC\,BF$ and $AC^2=BC\,FC$.
\end{theorem}
\begin{proof}
Use $A=0$, $B=u$, $C=v$, where $u\mathbin{\cdot}v=0$.
The projection parameter is
\[
 F=u+t(v-u),\qquad
 t=\frac{(-u)\mathbin{\cdot}(v-u)}{|v-u|^2}=\frac{c^2}{a^2}.
\]
Since $a^2=b^2+c^2$ and $b,c>0$, one has $0<t<1$.
Thus $BF=ta=c^2/a$ and $FC=(1-t)a=b^2/a$.
Orthogonality of $u,v$ now gives
\[
 |F|^2=(1-t)^2c^2+t^2b^2
 =\frac{b^4c^2+c^4b^2}{a^4}=\frac{b^2c^2}{a^2}.
\]
Taking its positive square root gives the altitude formula and its product
identity.  The projection itself is certified by
$(A-F)\mathbin{\cdot}(C-B)=0$, which follows from the definition of $t$.
\end{proof}

\subsection{Centers, concurrency, and extremal identities}

\begin{theorem}[Median concurrency and the centroid]
\label{tri:centroid}
The three medians of a noncollinear triangle meet at
$G=(A+B+C)/3$.  On each median, the vertex-to-centroid segment is twice
the centroid-to-midpoint segment.  This common point is unique.
\end{theorem}
\begin{proof}
For $M_A=(B+C)/2$,
$G=A+(2/3)(M_A-A)$, so $G$ lies strictly between $A,M_A$ and divides
the median in the asserted ratio.  Cycling gives all three incidences.
To check uniqueness, translate $A=0$, $B=u$, $C=v$.  A point on the first
median is $t(u+v)$, and one on the median from $B$ is
$u+r(v/2-u)$.  Linear independence of $u,v$ gives
$t=1-r$ and $t=r/2$, hence $r=2/3$, $t=1/3$ and the point is $G$.
Linear independence here follows directly from $[u,v]\ne0$, using
determinants to extract the two coefficients.
\end{proof}

\begin{theorem}[Leibniz's identity and the centroid minimizer]
\label{tri:leibniz}
For any point $X$ and $G=(A+B+C)/3$,
\[
 |X-A|^2+|X-B|^2+|X-C|^2
 =3|X-G|^2+\frac{a^2+b^2+c^2}{3}.
\]
Thus the centroid is the unique point minimizing the sum of squared
distances to the three vertices.
\end{theorem}
\begin{proof}
Let $p=A-G$, $q=B-G$, $r=C-G$, so $p+q+r=0$.
Expanding the three squares about $X-G$ makes all linear terms cancel,
leaving $3|X-G|^2+|p|^2+|q|^2+|r|^2$.
Expansion of $|p+q+r|^2=0$ gives
$2(p\mathbin{\cdot}q+q\mathbin{\cdot}r+r\mathbin{\cdot}p)
=-(|p|^2+|q|^2+|r|^2)$.
It follows by expanding the three pairwise squared distances that
$a^2+b^2+c^2=3(|p|^2+|q|^2+|r|^2)$.
The identity follows.  Its additional term $3|X-G|^2$ is nonnegative and
vanishes exactly when $X=G$, proving the constructive minimization claim.
\end{proof}

\begin{theorem}[Circumcenter, perpendicular bisectors, and circumradius]
\label{tri:circumcenter}
Every noncollinear triangle has a unique point $O$ equidistant from its
vertices; the three perpendicular bisectors meet there.  In positively
oriented coordinates $A=0$, $B=u$, $C=v$,
\[
 O=\frac{|v|^2Ju-|u|^2Jv}{2D},\qquad
 R=|O-A|=\frac{abc}{4\Delta}.
\]
Consequently the sine-law ratio is $a/\sin\alpha=b/\sin\beta=
c/\sin\gamma=2R$.
\end{theorem}
\begin{proof}
The equations $|O|^2=|O-u|^2=|O-v|^2$ are equivalent to
\[
 2O\mathbin{\cdot}u=|u|^2,\qquad
 2O\mathbin{\cdot}v=|v|^2.
\]
Substitution of the displayed $O$, using Lemma~\ref{tri:gram}, verifies
both.  If another point satisfies them, their difference has zero dot
product with both $u,v$; the nonzero determinant makes that difference
zero.  The first equation also says
$(O-u/2)\mathbin{\cdot}u=0$, the perpendicular bisector of $AB$;
similarly for the other sides.
Put $U=|u|^2=c^2$, $V=|v|^2=b^2$, $p=u\mathbin{\cdot}v$.
The squared numerator of $O$ is
\[
 |VJu-UJv|^2=V^2U+U^2V-2UVp
 =UV(U+V-2p)=a^2b^2c^2.
\]
Thus $R^2=a^2b^2c^2/(4D^2)$ and positivity selects
$R=abc/(2D)=abc/(4\Delta)$.
The final assertion follows from Theorem~\ref{tri:sine}.
\end{proof}

\begin{theorem}[Altitude concurrency and the orthocenter]
\label{tri:orthocenter}
The three altitudes of a noncollinear triangle meet in the unique point
$H=A+B+C-2O$, where $O$ is the circumcenter.
\end{theorem}
\begin{proof}
Translate $A=0$, $B=u$, $C=v$, so $H=u+v-2O$.
The circumcenter equations give
\[
 H\mathbin{\cdot}(v-u)
 =(u+v)\mathbin{\cdot}(v-u)-2O\mathbin{\cdot}(v-u)
 =(V-U)-(V-U)=0.
\]
This is the altitude through $A$.  Also
$(H-u)\mathbin{\cdot}v=(v-2O)\mathbin{\cdot}v=V-V=0$ and
$(H-v)\mathbin{\cdot}u=U-U=0$, proving the other two incidences.
The normals of the last two altitude equations are $v,u$, which are
linearly independent, so their intersection is unique by the same
determinant calculation as for the circumcenter.  Translating back gives
the stated coordinate-free affine formula.
\end{proof}

\begin{theorem}[Euler line]
\label{tri:euler-line}
The circumcenter $O$, centroid $G$, and orthocenter $H$ satisfy
$H=3G-2O$.  In particular they are collinear and $|H-G|=2|G-O|$.
If the points coincide, the affine identity remains valid; a distinct
``Euler line'' is asserted only when $O\ne H$.
\end{theorem}
\begin{proof}
The formulas $3G=A+B+C$ and $H=A+B+C-2O$ give the identity.
Then $G=O+(H-O)/3$, so collinearity and the ratio follow from
$H-G=2(G-O)$.  This proof includes the coincident case without attempting
to construct a line from two equal points.
\end{proof}

\begin{theorem}[Nine-point circle]
\label{tri:nine-point}
Let $N=(O+H)/2$.  The circle with center $N$ and radius $R/2$ contains
the three side midpoints, the three midpoints between $H$ and a vertex,
and the three altitude feet.  Some of these nine named points may
coincide.
\end{theorem}
\begin{proof}
Use the circumcenter as origin, and write the vertex vectors as $x,y,z$.
They have common length $R$, the orthocenter is $h=x+y+z$, and $N=h/2$.
The midpoint $M_A=(y+z)/2$ satisfies $M_A-N=-x/2$, while the midpoint
$K_A=(x+h)/2$ satisfies $K_A-N=x/2$.  Both therefore have distance
$R/2$ from $N$, and cycling gives six points.
Let $F_A$ be the projection of $x$ onto the supporting line of $yz$.
The points $M_A,F_A$ lie on that line.  Both $x$ and $h$ lie on the
altitude perpendicular to it, hence so does their midpoint $K_A$.
Therefore
$(F_A-M_A)\mathbin{\cdot}(F_A-K_A)=0$.
Since $N=(M_A+K_A)/2$, expansion rewrites this equality as
\[
 |F_A-N|^2=|M_A-K_A|^2/4=R^2/4.
\]
The positive distance root is $R/2$.  This proves the remaining three
incidences without requiring the feet to lie on the closed side segments.
\end{proof}

\begin{theorem}[Internal angle-bisector theorem and its converse]
\label{tri:bisector}
The internal bisector from $A$ meets $BC$ at
\[
 T=\frac{bB+cC}{b+c},\qquad \frac{BT}{TC}=\frac cb=\frac{AB}{AC}.
\]
Conversely an interior point of $BC$ with this ratio lies on that
internal bisector.
\end{theorem}
\begin{proof}
The unit directions from $A$ to $B,C$ are $e=(B-A)/c$ and
$f=(C-A)/b$.  Their sum $d=e+f$ is nonzero and lies inside the angle:
with positive orientation, $[e,d]=[d,f]=[e,f]>0$.
Also $e\mathbin{\cdot}d=f\mathbin{\cdot}d=1+e\mathbin{\cdot}f$.
After dividing by $|d|>0$, these equalities say that $d$ makes equal
angle pairs with $e,f$.  This certifies the internal bisector.  It is
unique as a ray: equality of those two cosines for a unit interior
direction $w$ gives $w\mathbin{\cdot}(e-f)=0$.  The nonzero perpendicular
vectors $e-f,e+f$ form a basis, so $w$ is a scalar multiple of $e+f$;
the interior side inequalities make that scalar positive.
The displayed point has
\[
 T-A=\frac{b(B-A)+c(C-A)}{b+c}
     =\frac{bc}{b+c}(e+f).
\]
It is on the desired ray, and $T=B+\frac c{b+c}(C-B)$ is strictly
between $B,C$.  Thus $BT/TC=c/b$.
Conversely an interior point has $T=B+t(C-B)$ with $0<t<1$.
The ratio condition $t/(1-t)=c/b$ uniquely gives $t=c/(b+c)$, hence
the displayed point and its bisector incidence.
\end{proof}

\begin{theorem}[Incenter, incircle, and the area--inradius identity]
\label{tri:incenter}
The three internal bisectors meet at
\[
 I=\frac{aA+bB+cC}{a+b+c}.
\]
This is the unique interior point equidistant from the three supporting
side lines.  Its common distance is
$r=2\Delta/(a+b+c)=\Delta/s>0$.
The perpendicular feet lie strictly on the side segments, so the circle
with center $I$ and radius $r$ is an incircle tangent to all three sides.
\end{theorem}
\begin{proof}
Put $P=a+b+c$.  Since
$I-A=\frac{bc}{P}((B-A)/c+(C-A)/b)$,
the preceding bisector calculation puts $I$ on the bisector from $A$;
cycling gives the other two.  Its barycentric coefficients $a/P,b/P,c/P$
are positive and sum to one, so it is strictly inside.
In the positively oriented coordinates $A=0$, $B=u$, $C=v$,
\[
 [C-B,I-B]=\frac{aD}{P},\quad
 [A-C,I-C]=\frac{bD}{P},\quad
 [B-A,I-A]=\frac{cD}{P}.
\]
These follow by bilinearity, since each determinant loses two of the
three barycentric terms.  Divide respectively by $a,b,c$ to get common
positive perpendicular distance $r=D/P$.
For uniqueness, any strictly interior $X$ with common distance $r'$ has
subtriangle doubled areas $ar'$, $br'$, $cr'$.  Their sum is $D$ by the
determinant identity
\[
 [C-B,X-B]+[A-C,X-C]+[B-A,X-A]=[B-A,C-A].
\]
Thus $r'=D/P$.  The first two prescribed determinants then determine
$X$ uniquely because their side directions are independent; $I$ already
has those values, so $X=I$.
Finally let $F$ be the foot on $BC$.  Its parameter from $B$ is
\[
 t=\frac{(I-B)\mathbin{\cdot}(C-B)}{a^2}
   =\frac{a+c-b}{2a}.
\]
For completeness, use
$I-B=(a(A-B)+c(C-B))/P$ and
$2(A-B)\mathbin{\cdot}(C-B)=c^2+a^2-b^2$.
The numerator simplifies to
$a((a+c)^2-b^2)/(2P)=a(a+c-b)/2$, proving the formula.
Strict triangle inequalities give $0<t<1$.  The same holds cyclically.
Each foot lies on the radius-perpendicular side line and at distance $r$
from $I$.  A point $F+t(C-B)$ on $BC$ has squared distance
$r^2+t^2a^2$ from $I$, so it lies on the circle only when $t=0$.
This certifies tangency and the incircle assertion, and the same
calculation applies to the other two sides.
\end{proof}

\begin{theorem}[Excenters and exradii]
\label{tri:excenters}
The point
\[
 I_A=\frac{-aA+bB+cC}{-a+b+c}
\]
has equal positive distances $r_A=\Delta/(s-a)$ from the three supporting
side lines.  It lies on the opposite side of $BC$ from $A$, and on the
same sides of $CA,AB$ as $B,C$, respectively.  The analogous formulas
give the other two excenters and radii.
\end{theorem}
\begin{proof}
The denominator $P_A=-a+b+c=2(s-a)$ is strictly positive.
Exactly the determinant calculation in the preceding proof, with the
coefficient of $A$ negated, gives signed doubled subareas
$-aD/P_A$, $bD/P_A$, $cD/P_A$.
After division by the corresponding side lengths, the signed distances
are $-D/P_A$, $D/P_A$, $D/P_A$.
Their absolute values are $D/P_A=\Delta/(s-a)>0$, proving the claim
and the three specified side relations.  Projection to each line gives
a tangent radius, so these data also construct its excircle.  This
statement concerns supporting lines; it does not silently claim all
three tangency points lie on the original closed segments.
\end{proof}

\begin{theorem}[Euler's inradius--circumradius identity and inequality]
\label{tri:euler-inequality}
For any noncollinear triangle,
\[
 |O-I|^2=R(R-2r),\qquad R\ge2r.
\]
Equality holds precisely for an equilateral triangle.
\end{theorem}
\begin{proof}
Use $O$ as origin, write the vertex vectors as $x,y,z$, and put
$P=a+b+c$.  They all have squared length $R^2$, while
$I=(ax+by+cz)/P$.  Expanding and collecting pairwise differences gives
\[
 |ax+by+cz|^2
 =P(a|x|^2+b|y|^2+c|z|^2)
   -ab|x-y|^2-bc|y-z|^2-ca|z-x|^2.
\]
The three pairwise distances here are $c,a,b$, respectively.  Thus
\[
 |I|^2=R^2-\frac{ab c^2+bc a^2+ca b^2}{P^2}
       =R^2-\frac{abc}{P}.
\]
The formulas $R=abc/(4\Delta)$ and $r=2\Delta/P$ give
$2Rr=abc/P$, proving the identity.
The left side is nonnegative and $R>0$, so $R-2r\ge0$.
If equality holds, $I=O$.  For each side, its midpoint is the foot from
the circumcenter, because the circumcenter lies on the perpendicular
bisector.  Since its distance to every side is now $r$, Pythagoras gives
\[
 a^2/4+r^2=b^2/4+r^2=c^2/4+r^2=R^2.
\]
Positive roots yield $a=b=c$.
Conversely for an equilateral triangle of side $\ell$,
$\Delta=\sqrt3\ell^2/4$ by Heron's formula, so
$r=\sqrt3\ell/6$ and $R=\ell/\sqrt3=2r$.
\end{proof}

\begin{theorem}[Viviani's theorem]
\label{tri:viviani}
For any point inside or on an equilateral triangle, the sum of its three
perpendicular distances to the side lines equals the altitude of the
triangle.
\end{theorem}
\begin{proof}
Let the common side be $\ell>0$, the altitude $h$, and the three
nonnegative signed distances of the point $X$ be $d_a,d_b,d_c$.
The determinant decomposition used in the incenter proof gives
\[
 2\Delta=\ell d_a+\ell d_b+\ell d_c.
\]
Also $2\Delta=\ell h$ from the perpendicular projection formula for
triangle area.  Cancel $\ell$ to obtain $d_a+d_b+d_c=h$.
Boundary points simply make one or more of the signed distances zero,
and require no separate limiting argument.
\end{proof}

\subsection{Arithmetic and Lean obligations for this chapter}

Every displayed construction has constant algebraic depth once its input
scalar codes are given.  The only roots introduced are nonnegative square
roots.  Denominators are products of positive side lengths, nonzero
determinants, $2$, $3$, or the explicitly positive quantities $s-a$, $s-b$,
$s-c$.  Hence the finite-code arithmetic layer can attach a nonzero
certificate to each division.  No minimization over an unbounded set is
used: even the centroid minimization theorem is proved by one identity.

The polynomial equalities should be exported as denominator-cleared ring
certificates.  Ordered consequences use only sums of squares, multiplication
by a certified positive scalar, and the rule that positive square roots of
equal nonnegative quantities are equal.  The two places that choose a
finite sign branch---the ``at least two acute angles'' conclusion and the
two solutions in the equilateral construction---use decidable comparison
of the exact scalar codes.  A Cauchy-real presentation may retain the
branch-free implications and use explicit apartness/side witnesses for
chosen branches; it must not import this decidability for arbitrary real
names.

A Lean development can first state these identities for a structure
\texttt{Vec2} with two scalar fields, define dot product, determinant and
quarter-turn by the formulas above, and check the denominator-cleared
identities with ring normalization.  The geometric layer then packages
the following data: a nonzero determinant for a triangle; positive square
roots for lengths; a unit-pair angle and its half-plane tag; an explicit
affine map for congruence; and a barycentric or projection parameter for
incidence.  The next stage proves the primitive-recursive coding and
certificate-soundness lemmas in the chosen \ipra{} embedding.  An ambient
Lean ordered-field proof alone is not that last stage, and the separation
must remain visible in the formalization status.
